% ============================================================================
\clearpage
\section*{v4.10.0-r04 Source-Only: The Vacuum Term with a Tachyon Branch, and the Bosons Created at the Big Bang}
\addcontentsline{toc}{section}{v4.10.0-r04 Source-Only: The Vacuum Term with a Tachyon Branch, and the Bosons Created at the Big Bang}
\label{sec:v4100-r04-tachyon-branch}
% ============================================================================

\paragraph{Purpose.}
Layer v4.10.0-r03 found that under the second reading no tachyon-sector mechanism gives a vacuum term after the big bang
(\S\ref{sec:v4100-r03-bang-g2}).  This layer records the check under the domain G1, a temporary author decision to admit the tachyon branch
as degrees of freedom of its own, two analyses of the deferred idea I3 under that decision (fixed weights and a dynamic branch gas, H95),
the withdrawal of the decision, an estimate of the bosons created at the bang, and the state of the comparison data of the open predictions.
The notes are \texttt{notes/\allowbreak research/\allowbreak v4\_10\_0\_r04/}.  Earlier layers are not rewritten.

\noindent\textbf{Source-only status.}
This v4.10.0-r04 layer is source-only.  It does not build a PDF, does not assign a new DOI,
and does not authorize a public release.  The v4.10.0 Version DOI
\href{https://doi.org/10.5281/zenodo.23991357}{10.5281/zenodo.23991357} and the Concept DOI
\href{https://doi.org/10.5281/zenodo.19210728}{10.5281/zenodo.19210728} are unchanged by it.

\subsection*{v4.10.0-r04 results}
\label{subsec:v4100-r04-results}

\paragraph{Theoremlet v4.10.0-r04.1 (domain G1 gives no tension either).}
Exploration (analytic).  The mixture of I3 (k182, k183) consists of a normal branch (\(\omega=m\cosh\eta\), \(k=m\sinh\eta\)) and a tachyonic
branch (\(\omega=m\sinh\eta\), \(k=m\cosh\eta\)); the tension comes from the tachyonic one.  The reduced theory of the second reading,
\((a-b)\tfrac12X-(a+b)V\), has \(\omega^2=k^2+\gamma m^2\) for every mode and no tachyonic branch.  For a massive scalar the averaged pressure
vanishes (dust).  Under G1, as under G2, the only vacuum term is a free multiple of \(V_{\min}\).  I3 needs the tachyonic branch as degrees of freedom
of its own, that is, a reversal of the decision of r84.

\paragraph{Decision v4.10.0-r04.2 (author, 6 November 2026: tachyon branch admitted, later withdrawn).}
First the tachyonic branch was admitted as degrees of freedom of its own, so that I3 becomes possible.  The known cost: quanta of negative energy, and
vacuum decay through the graviton with a frame cut-off \(\lesssim14\) MeV (r83, H76).  After r04.3 and r04.4 the author withdrew the decision: the
tachyon part is again only the second reading of the same field (r84).

\paragraph{Theoremlet v4.10.0-r04.3 (I3 with fixed weights is excluded).}
Exploration (analytic, with an estimate).  The M and T readings of the mixture have the same density and dilute alike.  With fixed weights the sum is
one fluid with constant \(w=-b/(3a)\), not dust plus \(\Lambda\).  To explain \(\Omega_\Lambda\approx0.69\) today one needs \(b/a\approx2\), and then the dark
component at recombination is about \(10^{-5}\) of the baryons instead of \(5.4\) (cosmological values from prior knowledge, not looked up).  In
addition, the mixture at common rapidity is not preserved under redshift: \(k_1/k_2=\tanh\eta\) would stay fixed while each \(k\) is fixed by \(\eta\).

\paragraph{Theoremlet v4.10.0-r04.4 (H95: a dynamic branch gas cannot do it either).}
Pre-registered (commit \texttt{9eb4e40}, not blind; the refutation was expected analytically).  Model: flat FRW with radiation and baryons, and a
freely streaming gas with equal numbers in both branches.  The normal branch has weight \(1\); the tachyonic branch has weight \(r\), negative
energy, and threshold at \(a_{\rm th}=k_0/m\).  Grid \(60\times60\times40\) in \((r,a_{\rm th},\kappa_1/a_{\rm th})\): of \(144\,000\) points, \(12\,780\) keep
the total density positive for \(a\in[10^{-8},1]\).  None has both a dark-to-baryon ratio in \([4.32,6.48]\) at \(z=1090\) and \(w_{\rm eff}\in[-0.8,-0.6]\)
today.  \textbf{Refuted.}
\emph{Noticed afterwards.}  \(w_0\) can be made arbitrarily negative (down to \(-151\)), but only with the threshold practically today
(\(a_{\rm th}=1.000001\)), and then the ratio at recombination is \(266\).

\paragraph{Theoremlet v4.10.0-r04.5 (bosons created at the bang; estimate).}
Exploration.  With \(\Delta\lesssim\ell_P\) (r03.3), bosons with matter weights receive \(n_0=0.404\) per mode up to \(k_{\max}\sim1/\Delta\gtrsim M_P\).  This
gives \(\rho\approx g\,n_0\,k_{\max}^4/(8\pi^2)=g\cdot0.0051\,k_{\max}^4\), equivalent to a temperature of \(\approx0.35\,k_{\max}\).  The radiation after the bang is
Planck-hot, non-thermal, and initially without fermions.  The number depends on the non-unitary range \(k\sim1/\Delta\) of r03.2 and is an order of
magnitude only.

\paragraph{Remark v4.10.0-r04.6 (state of the comparison data).}
For H15 (\(m_W\)) no world average published after 2026-10-21 was found, and for H42/H46 (\(\theta_{23}\)) no global analysis after 2026-10-24; all three
remain open.  \emph{Disclosure:} The PDG 2026 edition (May 2026) already gave \(m_W=80.3625\pm0.0077\) GeV before the pre-registration of H15, which
listed only PDG 2025 as known.  The value is \(0.13\sigma\) from the prediction and is not a test.

\paragraph{Proof and measurement sources.}
Script \texttt{k212\_\allowbreak zweiggas\_\allowbreak frw.py} with output, and the explorations and the proposal in
\texttt{notes/\allowbreak research/\allowbreak v4\_10\_0\_r04/}; pre-registration and result H95 in \texttt{notes/\allowbreak research/\allowbreak predictions/}.  \(\square\)

\paragraph{Consequence.}
The deferred idea I3 (\(\Lambda\) from the dark-matter calculation) has failed in both forms examined, with fixed weights and with a dynamic branch gas,
even with the tachyon branch admitted.  The author returned to the second reading alone, so \(\Lambda\) is an input.  The bosons created at the bang
form a Planck-hot, non-thermal start without fermions.  An observable trace would have to survive thermalization, for example a gravitational-wave
background, which first needs a rule for how tensor modes see the transition.  OP5 is not closed.

\subsection*{v4.10.0-r04 status in the three frames}
\label{subsec:v4100-r04-frames}

\begin{itemize}[leftmargin=2em]
\item r04.1--r04.4 concern how the gravity of frame M sees frame T.  Under the second reading T is a re-weighting of the same field (no own branch).
With an own branch, T would have quanta of negative energy; that was examined and withdrawn.
\item r04.5 holds in frame M after the bang.  With photon weights (frame P, \(a=b\)) no bosons are created.
\item No frame-free statement is made.
\end{itemize}

\subsection*{v4.10.0-r04 register entries}
\label{subsec:v4100-r04-register}

\begin{center}
\small
\begin{tabular}{@{}>{\raggedright\arraybackslash}p{0.34\textwidth}>{\raggedright\arraybackslash}p{0.58\textwidth}@{}}
\toprule
statement & status after v4.10.0-r04\\
\midrule
\(\Lambda\) under G1 (second reading) & no tension; free multiple of \(V_{\min}\)\\
tachyon branch as own degrees of freedom & admitted, then withdrawn (author)\\
I3 with fixed weights & excluded (constant \(w\), CMB)\\
I3 with dynamic branch gas & refuted (H95)\\
bosons created at the bang & \(\rho\approx g\cdot0.0051\,k_{\max}^4\), non-thermal, no fermions (estimate)\\
H15, H42, H46 & open; PDG 2026 \(m_W\) disclosed as known before H15\\
\bottomrule
\end{tabular}
\end{center}

\subsection*{v4.10.0-r04 ledger and guard update}
\label{subsec:v4100-r04-ledger-guard}

\[
R_{\rm chk}^{v4100\text{-}r04}=(1,1,0)
\]
for (I3 excluded in both forms; decision withdrawn; \(\Lambda\) remains an input).

The layer spends no guard:
\[
R_{\rm glob}^{(v4100\text{-}r04,+)}=(0,1,1,1,1),
\qquad
\Gamma_{\rm guard}^{v4100\text{-}r04}=(1,1,1,1,1),
\qquad
\beta_4=1/256\ \text{unspent}.
\]
It does not imply exact finite \(C5\) zero closure, exact all-readout terminality,
\(C6/T2\), a global photon-clock/RH theorem, or global all-readout closure.
